\chapter{คู่มือการติดตั้งและตั้งค่าสภาพแวดล้อมการทำงาน}
\label{app:environment_setup}
\index{การติดตั้งสภาพแวดล้อม}
\index{Google Colab}

ภาคผนวกนี้รวบรวมขั้นตอนการติดตั้งและกำหนดค่าสภาพแวดล้อมการทำงาน (Development Environment Setup) สำหรับการพัฒนาและทดสอบระบบปัญญาประดิษฐ์วินิจฉัยโรคใบทุเรียน ทั้งบนระบบคลาวด์ Google Colab และบนเครื่องคอมพิวเตอร์ส่วนบุคคล

\section{การเตรียมสภาพแวดล้อมบน Google Colab (Cloud GPU)}
Google Colab เป็นสภาพแวดล้อมที่เหมาะสมที่สุดสำหรับผู้เริ่มต้นที่ไม่มีการ์ดจอ GPU โดยมีขั้นตอนการตั้งค่าดังนี้:

\begin{enumerate}[leftmargin=1.8em, itemsep=4pt]
    \item \textbf{การสร้างสมุดงานและเลือกการ์ดจอ GPU:}
    \begin{itemize}[leftmargin=1.2em, itemsep=2pt]
        \item เข้าสู่เว็บไซต์ \texttt{https://colab.research.google.com} และลงชื่อเข้าใช้ด้วยบัญชี Google
        \item ไปที่เมนู \textit{Runtime $\rightarrow$ Change runtime type}
        \item ในช่อง \textit{Hardware accelerator} เลือก \textbf{T4 GPU} และกดปุ่ม \textit{Save}
    \end{itemize}
    
    \item \textbf{การติดตั้งและอัปเดตแพ็กเกจที่จำเป็น:}
    พิมพ์คำสั่งต่อไปนี้ในเซลล์แรกของสมุดงาน:
\begin{lstlisting}[language=bash]
!pip install --upgrade torch torchvision torchaudio
!pip install onnx onnxruntime opencv-python matplotlib pillow torchinfo
\end{lstlisting}

    \item \textbf{การตรวจสอบสถานะการเชื่อมต่อ GPU:}
\begin{lstlisting}[language=Python]
import torch
print(f"[*] PyTorch Version: {torch.__version__}")
print(f"[*] CUDA Available: {torch.cuda.is_available()}")
if torch.cuda.is_available():
    print(f"[*] GPU Device Name: {torch.cuda.get_device_name(0)}")
\end{lstlisting}
\end{enumerate}

\section{การติดตั้งสภาพแวดล้อมบนเครื่องคอมพิวเตอร์ส่วนบุคคล (Local Machine)}
สำหรับผู้พัฒนาที่ต้องการรันระบบบนเครื่องคอมพิวเตอร์ส่วนบุคคล (รองรับทั้ง Windows, macOS และ Linux):

\begin{enumerate}[leftmargin=1.8em, itemsep=4pt]
    \item \textbf{การสร้างสภาพแวดล้อมเสมือน (Python Virtual Environment):}
\begin{lstlisting}[language=bash]
# สร้างโฟลเดอร์โครงการและสภาพแวดล้อมเสมือน
python3 -m venv venv_durian

# การเปิดใช้งานสภาพแวดล้อม (macOS/Linux)
source venv_durian/bin/activate

# การเปิดใช้งานสภาพแวดล้อม (Windows PowerShell)
.\venv_durian\Scripts\Activate.ps1
\end{lstlisting}

    \item \textbf{การติดตั้งไลบรารีจากไฟล์ requirements.txt:}
\begin{lstlisting}[language=bash]
pip install -r requirements.txt
\end{lstlisting}
\end{enumerate}

\section{การติดตั้งชุดพัฒนา Flutter สำหรับการสร้างแอปพลิเคชันมือถือ}
\index{Flutter SDK}
การสร้างแอปพลิเคชันสมาร์ทโฟนสำหรับติดตั้งบนระบบปฏิบัติการ Android และ iOS มีขั้นตอนดังนี้:
\begin{enumerate}[leftmargin=1.8em, itemsep=3pt]
    \item ดาวน์โหลด Flutter SDK เวอร์ชันล่าสุดจากเว็บไซต์ทางการ \texttt{https://flutter.dev}
    \item แตกไฟล์และเพิ่มพาธของโฟลเดอร์ \texttt{flutter/bin} เข้าสู่ System Environment Variables (PATH)
    \item เปิดเทอร์มินัลและรันคำสั่งตรวจสอบความพร้อม: \texttt{flutter doctor}
    \item ติดตั้งปลั๊กอินสำหรับ ONNX Runtime ในไฟล์ \texttt{pubspec.yaml}:
\begin{lstlisting}[language=bash]
dependencies:
  flutter:
    sdk: flutter
  camera: ^0.10.5
  onnxruntime: ^1.2.0
  path_provider: ^2.1.2
\end{lstlisting}
\end{enumerate}
